\chapter{Podręcznik użytkownika}

W tym rozdziale opisano sposób kompilacji programu, obsługiwane argumenty programu oraz konfigurację parametrów dotyczących automatów oraz generatorów.

\section{Uruchomienie programu}

Projekt kompilowany jest przy użyciu skryptu \texttt{build.sh}, który umożliwia:

\begin{itemize}
    \item Standardowe zbudowanie pliku wykonywalnego \path{build/run} przy pomocy komendy: \texttt{./build.sh}
    \item Zbudowanie pliku wykonywalnego \path{build-debug/run} z dodatkowymi flagami debugowymi przy pomocy komendy: \texttt{./build.sh -d}
    \item Zbudowanie pliku wykonywalnego, w którym zapytania o równoważność są sprawdzane heurystycznie \path{build-random/run} przy pomocy komendy: \\ \texttt{./build.sh -randomEQ}
    \item Zbudowanie i uruchomienie wszystkich testów jednostkowych używając komendy: \texttt{./build.sh test} (lub \texttt{./build.sh -d test})
    \item Zbudowanie i uruchomienie pojedynczego testu jednostkowego \texttt{testName} używając komendy: \texttt{./build.sh test -t testName} (lub \texttt{./build.sh -d test -t testName})
\end{itemize}

Program można uruchomić przy użyciu następujących poleceń:
\begin{verbatim}
    ./run <generator> <seed> <numOfTests>
    ./run bench <scenario>
    ./run custom
    ./run --help
    
Generatory:
    random        uruchomienie testów z losowym generatorem
    cda           uruchomienie testów z generatorem CDA
    sevpa         uruchomienie testów z generatorem SeVpa
    mevpa         uruchomienie testów z generatorem MeVpa
    ecda          uruchomienie testów z generatorem eCDA
    commutative   uruchomienie testów z generatorem commutative
    cancel        uruchomienie testów z generatorem cancellation
    idempotency   uruchomienie testów z generatorem idempotency
    
Pozostałe komendy:
    bench      uruchomienie wybranego scenariusza benchmarkowego
    custom     wykonanie algorytmu dla przykładowego automatu
               z pliku `helper/ExampleRun.cpp`

Scenariusze:
    Nazwy scenariuszy zdefiniowane są w pliku:
    `benchmark/scenario/scenarios.hpp`

Argumenty:
    seed         ziarno generatora liczb losowych
    numOfTests   liczba testów do wykonania
    scenario     nazwa scenariusza benchmarkowego

Przykłady uruchomień:
    ./run random 42 100
    ./run cda 123 50
    ./run bench base-scenario
    ./run custom
\end{verbatim}

\section{Konfiguracja programu}

Program zawiera dwa istotne pliki konfiguracyjne:
\begin{enumerate}
    \item \texttt{utils/Constants.hpp} zawiera m.in. definicje parametrów: \texttt{MaxNumOf\-Automa\-ton\-States}, \texttt{MaxNumOfLetters}
    oraz \texttt{MaxNumOfStackSymbols}, które określają maksymalny rozmiar automatów przyjmowanych przez algorytm. 
    Zawiera również parametry \texttt{MaxNumOfWellMatchedWordsForSrsCheck} oraz \texttt{MaxNumOf\-Rev\-ealed\-Well\-MatchedWordsForSrsCheck} określające maksymalną ilość dobrze zbalansowancyh słów do sprawdzania
    zgodności z SRS oraz maksymalną ilość rozważanych słów, które potencjalnie mogą zostać wybrane do sprawdzenia zgodności.

    \item \texttt{TesterParameters.hpp} zawiera definicję struktury \texttt{TesterParameters} służącej do parametryzowania testu oraz obiekty:
    \texttt{<nazwa\_generatora>Test\-Param\-eters}, które są wykorzystywane podczas wykonywania programu dla odpowiedniego generatora.
\end{enumerate}

Zmiana powyższych parametrów konfiguracyjnych wymaga ponownej kompilacji programu.